\exercisesection{\S~10}

\begin{enumerate}
\def\labelenumi{\arabic{enumi})}
\item
  令 \(\mathfrak{s} = \mathfrak{sl}(2, k)\)，\(\mathfrak{g} = \mathfrak{sl}(3, k)\)。把 \(\mathfrak{s}\) 与 \(\mathfrak{g}\) 的一个子代数相识别（借助 \(\mathfrak{s}\) 的 3 次不可约表示）。证明 \(\mathfrak{g}\) 的包含 \(\mathfrak{s}\) 且在 \(\operatorname{ad}_{\mathfrak{g}}\mathfrak{s}\) 下稳定的子空间都等于 \(\mathfrak{s}\) 或 \(\mathfrak{g}\)；由此推出 \(\mathfrak{s}\) 在 \(\mathfrak{g}\) 的异于 \(\mathfrak{g}\) 的子代数中是极大的。
\item
  令 \(m = \frac{1}{2}(\dim(\mathfrak{g}) + \operatorname{rk}(\mathfrak{g}))\)。g 的每个可解子代数的维数都 \(\leq m\)；若其维数为 m，则它是 Borel 子代数。（化归到代数闭的情形，并使用定理2。）
\item
  假设 \(k\) 是 \(\mathbf{R}, \mathbf{C}\)，或非离散完备超度量域。赋予格拉斯曼流形 \(\mathbf{G}(\mathfrak{g})\)（由 \(\mathfrak{g}\) 的向量子空间构成）其在 \(k\) 上的自然解析流形结构（《微分与分析流形：结果》，5.2.6）。考察 \(\mathbf{G}(\mathfrak{g})\) 中由下列对象组成的子集：
\end{enumerate}

\begin{enumerate}
\def\labelenumi{(\roman{enumi})}
\item
  子代数
\item
  可解子代数
\item
  幂零子代数
\item
  由幂零元素组成的子代数
\item
  Borel 子代数。
\end{enumerate}

证明这些子集都是闭的（对 (v)，使用习题2）。由此推出，当 \(k\) 局部紧时，这些子集都是紧的。

用例子表明，\(\mathbf{G}(\mathfrak{g})\) 中由下列对象组成的子集：

\begin{enumerate}
\def\labelenumi{(\roman{enumi})}
\setcounter{enumi}{5}
\item
  嘉当子代数
\item
  在 \(\mathfrak{g}\) 中约化的子代数
\item
  半单子代数
\item
  可分解子代数
\end{enumerate}

即使当 \(k = \mathbf{C}\) 时，也未必是闭的。

\begin{enumerate}
\def\labelenumi{\arabic{enumi})}
\setcounter{enumi}{3}
\tightlist
\item
  假设 \(k = \mathbf{C}\)。令 \(G = \operatorname{Int}(\mathfrak{g}) = \operatorname{Aut}_0(\mathfrak{g})\)，并令 \(B\) 为 \(G\) 的积分子群，其李代数 \(\mathfrak{b}\) 是 \(\mathfrak{g}\) 的 Borel 子代数。证明 \(B\) 是 \(\mathfrak{b}\) 在 \(G\) 中的正规化子（使用§3的习题4和§5的习题11）。借助习题3推出 \(G / B\) 是紧的。
\end{enumerate}

¶5) 假设 g 可分裂。若 h 是 g 的分裂嘉当子代数，则以 \(\mathrm{E}(\mathfrak{h})\) 表示 \(\mathrm{Aut}_{e}(\mathfrak{g})\) 中由诸 \(e^{\operatorname{ad} x}, x \in \mathfrak{g}^{\alpha}(\mathfrak{h}), \alpha \in \mathrm{R}(\mathfrak{g}, \mathfrak{h})\) 生成的子群，参见第VII章§3第2节。

\begin{enumerate}
\def\labelenumi{\alph{enumi})}
\item
  令 \(\mathfrak{b}\) 为 \(\mathfrak{g}\) 的包含 \(\mathfrak{h}\) 的 Borel 子代数，令 \(\mathfrak{h}_1\) 为 \(\mathfrak{b}\) 的嘉当子代数。证明 \(\mathfrak{h}_1\) 与 \(\mathfrak{h}\) 经 \(\mathrm{E}(\mathfrak{h})\) 的一个元素共轭。由此推出 \(\mathrm{E}(\mathfrak{h}) = \mathrm{E}(\mathfrak{h}_1)\)。
\item
  令 \(\mathfrak{h}'\) 为 \(\mathfrak{g}\) 的分裂嘉当子代数。证明 \(\mathrm{E}(\mathfrak{g}) = \mathrm{E}(\mathfrak{h}')\)。（若 \(\mathfrak{b}'\) 是包含 \(\mathfrak{h}'\) 的 Borel 子代数，选取 \(\mathfrak{h}_1\) 为 \(\mathfrak{b} \cap \mathfrak{b}'\) 的一个嘉当子代数，并应用 \(a\) ) 证明 \(\mathrm{E}(\mathfrak{h}) = \mathrm{E}(\mathfrak{h}_1) = \mathrm{E}(\mathfrak{h}')\)。）
\item
  令 \(x\) 为 \(\mathfrak{g}\) 的幂零元素。证明 \(e^{\mathrm{ad}x} \in \mathrm{E}(\mathfrak{h})\)。（利用 b) 和定理1的推论2，化归到 \(x \in [\mathfrak{b}, \mathfrak{b}]\) 的情形。）由此推出 \(\mathrm{E}(\mathfrak{h}) = \mathrm{Aut}_e(\mathfrak{g})\)。
\end{enumerate}

\begin{enumerate}
\def\labelenumi{\arabic{enumi})}
\setcounter{enumi}{5}
\tightlist
\item
  \begin{enumerate}
  \def\labelenumii{\alph{enumii})}
  \tightlist
  \item
    证明下列性质等价：
  \end{enumerate}
\end{enumerate}

\begin{enumerate}
\def\labelenumi{(\roman{enumi})}
\item
  g 没有 \(\neq 0\) 的幂零元素。
\item
  g 没有异于 g 的抛物子代数。
\end{enumerate}

（使用定理1的推论2。）

称这样的代数为各向异性的。

\begin{enumerate}
\def\labelenumi{\alph{enumi})}
\setcounter{enumi}{1}
\tightlist
\item
  令 \(\mathfrak{p}\) 为 \(\mathfrak{g}\) 的极小抛物子代数，\(\mathfrak{r}\) 为 \(\mathfrak{p}\) 的根基，\(\mathfrak{s} = \mathfrak{p} / \mathfrak{r}\)。证明 \(\mathfrak{s}\) 是各向异性的。（注意，若 \(\mathfrak{q}\) 是 \(\mathfrak{s}\) 的抛物子代数，则 \(\mathfrak{q}\) 在 \(\mathfrak{p}\) 中的逆像是 \(\mathfrak{g}\) 的抛物子代数，参见§3习题5 a)。）
\end{enumerate}

¶ 7) a) 证明 k 的下列性质等价：

\begin{enumerate}
\def\labelenumi{(\roman{enumi})}
\item
  每个各向异性的半单 \(k\)-李代数都化为 0。
\item
  每个半单 k-李代数都具有 Borel 子代数。
\end{enumerate}

（为证 (i) \(\Longrightarrow\) (ii)，使用习题6。）

\begin{enumerate}
\def\labelenumi{\alph{enumi})}
\setcounter{enumi}{1}
\tightlist
\item
  证明 (i) 和 (ii) 蕴含 \(^{19}\)：
\end{enumerate}

\begin{enumerate}
\def\labelenumi{(\roman{enumi})}
\setcounter{enumi}{2}
\tightlist
\item
  每个作为域的有限维 \(k\)-代数都是交换的。（或者再说一遍：\(k\) 的每个代数扩张的 Brauer 群都化为 0。）
\end{enumerate}

（使用该代数中迹为零的元素所成的李代数。）

\begin{enumerate}
\def\labelenumi{\alph{enumi})}
\setcounter{enumi}{2}
\tightlist
\item
  证明 (i) 和 (ii) 由下列性质蕴含：
\end{enumerate}

\begin{enumerate}
\def\labelenumi{(\roman{enumi})}
\setcounter{enumi}{3}
\tightlist
\item
  对任意由齐次多项式 \(f_{\alpha} \in k[(X_i)_{i \in \mathrm{I}}]\)（次数 \(\geq 1\)）组成且满足 \(\sum_{\alpha} \deg f_{\alpha} < \operatorname{Card}(\mathrm{I})\) 的有限族，存在不全为零的元素 \(x_i \in k\)，使得 \(f_{\alpha}((x_i)_{i \in \mathrm{I}}) = 0\) 对所有 \(\alpha\) 成立。
\end{enumerate}

（使用§8的命题5。）

